%% kapa: exam=1 round=19 year=2008 period=2 form=A subject=economics subject_ko=경제원론 seq=24 number=24
%% source: Q-Net 감정평가사 기출문제 (cst00309) artlSeq=1001962 file=extracted/1008/1001962/2008년 제19회 감정평가사 1차시험 기출문제/2교시.hwp
%% conversion: kordoc 4.24.51 -> tools/gen1.py (auto); review=proofread-sonnet; flags=none
\begin{question}{24}
완전경쟁적인 생산물시장에 참여하고 있는 어떤 기업은 생산요소시장에서 수요를 독점하고 있다. 이 기업이 이윤극대화를 위해 생산요소의 투입량을 어떻게 결정하는가? (단, 생산요소시장에서 요소의 공급곡선은 우상향한다고 가정한다.)
\begin{choices}
  \item 한계생산물가치곡선과 한계요소비용곡선이 만나는 점에서 요소 투입량을 결정한다.
  \item 한계수입곡선과 요소의 공급곡선이 만나는 점에서 요소 투입량을 결정한다.
  \item 한계수입곡선과 한계요소비용곡선이 만나는 점에서 요소 투입량을 결정한다.
  \item 한계수입생산곡선과 요소의 공급곡선이 만나는 점에서 요소 투입량을 결정한다.
  \item 한계수입생산곡선과 평균요소비용곡선이 만나는 점에서 요소 투입량을 결정한다.
\end{choices}
\end{question}
\answer{1}
